\section{Arquitectura del sistema: agente único, sandbox de máquina virtual y núcleo de Coding}

\subsection{Arquitectura central: Model + Sandbox + Tooling}

Cada sesión corresponde a una máquina virtual independiente, lo que genera aislamiento de recursos y de permisos. El modelo no «imagina» directamente el resultado de la ejecución, sino que en un entorno observable invoca comandos, escribe código, lee la retroalimentación y luego planea el siguiente paso. Esta es la clave para elevar al LLM de «contestador» a «ente ejecutor».

\begin{importantbox}{Por qué la máquina virtual es una decisión estructural}
\begin{itemize}
    \item Seguridad: la ejecución de la tarea queda aislada del equipo del usuario
    \item Continuidad: las tareas largas pueden seguir ejecutándose en segundo plano
    \item Observabilidad: los registros completos sirven para depurar, evaluar y reproducir
\end{itemize}
\end{importantbox}

\begin{figure}[H]
\centering
\includegraphics[width=0.92\textwidth]{waveform-full.png}
\caption{Forma de onda de la actividad de voz durante toda la entrevista (intensidad de audio aproximada)\protect\footnotemark}
\end{figure}
\footnotetext{Fuente: extraído del audio local de la entrevista `audio.m4a`. Intervalo de tiempo correspondiente: 00:00:00--03:27:00.}

\subsection{Prioridad del agente único frente a la orquestación Multi-Agent}

Ji Yichao (季逸超) manifiesta explícitamente sus reservas ante la tendencia actual de «forzar la división del trabajo entre múltiples agentes». La razón no es que Multi-Agent sea siempre incorrecto, sino que, dado el nivel actual de capacidad de los modelos y de madurez de la ingeniería, el costo de sincronizar estados y de la propagación de errores en la colaboración entre múltiples agentes supera el beneficio.

\begin{warningbox}{Los modos de falla más comunes en las etapas tempranas de Multi-Agent}
\begin{enumerate}
    \item La tarea se divide en fragmentos demasiado pequeños, y el costo de comunicación consume el beneficio de la ejecución
    \item La desincronización de estados provoca trabajo duplicado u operaciones en conflicto
    \item Es difícil localizar la falla y los límites de responsabilidad no están claros
\end{enumerate}
\end{warningbox}

\subsection{Coding como espacio de acción unificado}

En la entrevista, «Coding es el alma» no es una consigna. El valor esencial de Coding es unificar el espacio de acción: ya sea el procesamiento de datos, la generación de páginas web, las pruebas automatizadas o la orquestación de herramientas, todo puede reducirse a código y ejecución de comandos. De esta forma, la complejidad del sistema no crece linealmente con el número de herramientas.

\begin{knowledgebox}{Beneficios de un espacio de acción unificado}
\begin{itemize}
    \item Reduce la dificultad de aprender la política: el modelo solo necesita dominar un conjunto de habilidades muy reutilizable
    \item Reduce el costo de mantenimiento del producto: los cambios en las interfaces de las herramientas quedan absorbidos por la capa de código
    \item Aumenta la interpretabilidad: las fallas pueden reproducirse mediante registros y la reproducción de scripts
\end{itemize}
\end{knowledgebox}

\subsection{Diseño de herramientas: «quitar herramientas» en lugar de «agregar herramientas»}

A diferencia de muchos equipos de agentes que tienden a agregar herramientas continuamente, Manus enfatiza más la contención. Porque cuantas más herramientas hay, mayor es el espacio de selección de la política y mayor la probabilidad de una invocación errónea. Un conjunto pequeño de herramientas generales de alta calidad suele ser más estable que una gran cantidad de herramientas especializadas.

\begin{importantbox}{Tres criterios para gobernar las herramientas}
\begin{enumerate}
    \item Por defecto, intentar primero con la cadena de herramientas generales (Code Execution, Shell, Browser)
    \item Solo consolidar una herramienta especializada cuando la tasa de reutilización sea suficientemente alta
    \item En cada iteración de versión, evaluar la «lista de herramientas que pueden eliminarse»
\end{enumerate}
\end{importantbox}

\begin{figure}[H]
\centering
\includegraphics[width=0.92\textwidth]{waveform-mid.png}
\caption{Forma de onda de voz durante la etapa de discusión de la arquitectura (interacción de alta densidad en el tramo intermedio)\protect\footnotemark}
\end{figure}
\footnotetext{Fuente: extraído del audio local de la entrevista. Intervalo de tiempo correspondiente: 01:20:00--02:05:00.}

\subsection{Resumen de la sección}

La elección de arquitectura de Manus puede resumirse en una frase: lograr el máximo de tareas con la mínima estructura. El agente único, la ejecución en sandbox, el núcleo de Coding y la contención en las herramientas sirven en conjunto a la función objetivo de «completar tareas complejas de manera estable».

